% TOP_PROBLEM: {"id": 3212, "title": "Graceful Tree Conjecture", "category_id": 3, "category": {"id": 3, "name": "graph_theory", "display_name": "Graph Theory", "description": "Problems involving graphs, networks, and their properties.", "slug": "graph-theory", "order_index": 3, "created_at": "2026-07-31T15:26:25.671Z"}, "status": "open", "problem_number": "OPG-439", "set_id": 12, "statusQualification": "Injectivity of vertex labels is explicit, correcting the omission noted in the source discussion."}
% FIRST_POSTED: 2026-10-01
% ATTEMPT_STATUS: unresolved
% MODEL: GPT 6 Astra Ultra
% UnsolvedMath ID 3212; source catalogue code OPG-439
% !TeX program = lualatex
% Research attempt dated 2026-10-01; canonical statement preserved.
\documentclass[11pt,a4paper]{article}
\usepackage[margin=25mm]{geometry}
\usepackage{fontspec}
\setmainfont{lmroman10-regular.otf}[BoldFont=lmroman10-bold.otf,
 ItalicFont=lmroman10-italic.otf,BoldItalicFont=lmroman10-bolditalic.otf]
\usepackage{amsmath,amssymb,mathtools}
\usepackage{unicode-math}
\setmathfont{latinmodern-math.otf}
\AtBeginDocument{\let\setminus\smallsetminus}
\usepackage{xurl,hyperref}
\hypersetup{colorlinks=true,urlcolor=blue}
\setcounter{secnumdepth}{0}
\setlength{\parindent}{0pt}
\setlength{\parskip}{5pt}
\setlength{\emergencystretch}{3em}
\begin{document}
\section*{Graceful Tree Conjecture}\label{3212}
{\small UnsolvedMath ID 3212 · OPG-439 · Graph Theory · OpenGarden}
\subsection{Definitions and mathematical statement}
Let $T=(V,E)$ be a finite tree, meaning
a nonempty simple connected graph with
no cycle. Write $m=|E|$, so $|V|=m+1$.
A graceful labelling is an injective
function
\[
 f:V\longrightarrow\{0,1,\ldots,m\}
\]
for which the induced edge labels
$|f(u)-f(v)|$, $uv\in E$, are pairwise
distinct. Since there are $m$ edges
and each edge difference lies in
$\{1,\ldots,m\}$, this is equivalent
to requiring
\[
 \{|f(u)-f(v)|:uv\in E\}=\{1,\ldots,m\}.
\]
For a tree the vertex labelling is
necessarily a bijection, since domain
and codomain have the same size.

The Graceful Tree Conjecture states
that every finite tree admits a
graceful labelling. Both the vertex
labels and the absolute edge differences
must satisfy the specified conditions
simultaneously. Repeated vertex labels
are forbidden even if the edge
differences happen to be distinct.

The one-vertex tree, with $m=0$,
has the label zero and no edges,
so it satisfies the definition
vacuously. In general, no root,
orientation, or ordering of the
vertices is prescribed; the labelling
may be chosen freely to suit the
tree. Enlarging the label interval
or allowing fewer than $m$ distinct
edge differences would give weaker
labelling problems. The target here
requires exactly the interval from
zero to $m$ and every positive
difference in that interval once.
\subsection{Short English statement}
Can every tree be numbered from zero to its number of edges so the differences across its edges are exactly all positive labels, without repetition?
\subsection{Research attempt}
\paragraph{Scope and method.}
This research attempt by GPT-6 Astra Ultra is dated 1 October 2026.
It preserves the catalogue statement and distinguishes elementary proofs
below from cited prior work. No novelty claim or formal proof-assistant
verification is made. The final paragraph states the precise remaining scope.
\paragraph{Literature and constructive target.}
The catalogue source attributes graceful tree labelling to Kotzig, Ringel
and Rosa. The primary paper [R1] studies additional restricted tree
families. Here we construct explicit labellings for paths and all
double-stars, prove a cyclic decomposition consequence, and test the
most immediate leaf-extension strategy. The constructions are elementary
known cases, not evidence that all tree shapes have been covered.

\paragraph{Paths and double-stars.}
For a path with $m$ edges, assign labels along it as
$0,m,1,m-1,2,m-2,\ldots$, stopping when all $m+1$ labels are used.
Successive absolute differences are $m,m-1,\ldots,1$, so the labelling
is graceful. The single vertex with $m=0$ is labelled zero.

For a double-star, let its adjacent centres $u,v$ have respectively
$a,b$ additional leaf neighbours; thus $m=a+b+1$. Assign
\[
 f(u)=0,\qquad f(v)=b+1,
\]
give the $b$ leaves at $v$ labels $1,\ldots,b$, and give the $a$ leaves
at $u$ labels $b+2,\ldots,m$. Every integer from zero to $m$ is used once.
The differences at $v$ are $1,\ldots,b$; the centre edge has difference
$b+1$; and the differences at $u$ are $b+2,\ldots,m$. These three sets
partition the required edge labels. Empty leaf sets cause no problem,
so stars and the one-edge tree are included. This proves all trees of
diameter at most three, since deleting the central edge of such a tree
leaves two stars (with a star or single vertex as the degenerate cases).

\paragraph{A cyclic consequence with an exact count.}
Suppose a tree with $m\ge1$ edges has a graceful labelling. Regard its
labels as residues modulo $N=2m+1$ and translate all labels by each
$t\in\mathbb Z/N\mathbb Z$. The resulting $N$ embedded copies partition
the edges of $K_N$. Indeed unordered pairs have difference classes
$\{d,-d\}$ with $1\le d\le m$. For each class the original tree has
exactly one edge. Translating that edge produces every pair in its class
exactly once: a nonzero translation fixing an unordered pair would swap
its endpoints and have order two, impossible in a group of odd order.
Different classes cannot overlap. The total $Nm=\binom N2$ also matches.
This proves a consequence of graceful labelling; a decomposition obtained
by a different method need not provide the particular labels required here.

\paragraph{Why appending one new label is not an induction.}
Start with a graceful labelling using $0,\ldots,m$ and attach a leaf at
a vertex labelled $x$. If old labels are frozen, the only unused label
in the new interval is $m+1$. The only missing edge difference is also
$m+1$, so the new edge works exactly when $x=0$. Complementing all old
labels interchanges zero and $m$ but does not put an arbitrary prescribed
vertex at zero. For example, the path labelled $0,2,1$ is graceful, but
attaching a leaf at the endpoint labelled one and assigning it three
repeats difference two. The new path does have a graceful labelling,
$0,3,1,2$, after a global change. Thus this is an obstruction to a frozen
labelling algorithm, not a counterexample tree.

\paragraph{Verification and final status.}
The root batch script checks path labels through forty edges and every
double-star with $0\le a,b\le12$, comparing both the vertex-label and
edge-difference sets exactly. It also checks the complete-graph cyclic
decompositions for the first twelve paths and the failed leaf extension.
Those computations test formulas; the proofs above establish the infinite
families. \textbf{UNRESOLVED.} Arbitrary trees require a mechanism for
coordinated relabelling after adding branches. Assuming that a graceful
labelling can always put a prescribed attachment vertex at zero would
introduce an additional unproved assertion rather than complete induction.

\subsection{Sources}
\begin{itemize}
\item[S1] ULAM.AI and the UnsolvedMath Contributors, supplied problems.json, record 3212 (OPG-439), OpenGarden collection. Catalogue identity and source transcription; CC BY 4.0.
\url{https://huggingface.co/datasets/ulamai/UnsolvedMath}.
\item[S2] Open Problem Garden, Graceful Tree Conjecture. Original problem and discussion; the mathematical formulation below is expanded from the supplied source record.
\url{http://www.openproblemgarden.org/op/graceful_tree_conjecture}.
\item[R1] Matthew Superdock, \emph{The Graceful Tree Conjecture:
a class of graceful diameter-6 trees}, arXiv:1403.1564.
Primary source checked 1 October 2026.
\url{https://arxiv.org/abs/1403.1564}.
\end{itemize}
\end{document}
